\section{Additional Results}
\label{sec:additional-results}
% 【中文】本文件已含 v1.7 时期中文对照注释（`% <chinese>` 行内格式，功能等价；2026-11-18 中文对照令核验保留，未改以守判据③纯新增）。

% 附录 B：扩展实验结果（参考 SWE-bench Appendix C "Additional Results"）
% 附录 B：扩展实验结果（参考 SWE-bench Appendix C "Additional Results"）

\subsection{Full Per-Domain Results}
\label{sec:full-per-domain}
% \subsection{完整的各领域结果}

{\hbadness=10000\sloppy Table~\ref{tab:full-domain-results} reports the per-domain breakdown from the 3-run M0 evaluation (3 runs across all 13 domains, DeepSeek-V4-Pro, Cline v3.0.51, \texttt{thinking=low}, 16 workers, 2400s timeout). The overall Pass@1 is \MzeroPass{}$\pm$\MzeroStd{} (3-run mean$\pm$std; r1=105/482, r2=96/486, r3=103/490 decided; 304 PASS / 1154 FAIL / 42 ERROR-or-ABSENT across 1500 runs). Per-domain breakdown for the 3-run evaluation is pending and will be reported in the camera-ready version; the table below is retained as a placeholder for the per-domain matrix.\par}
% 表~\ref{tab:full-domain-results} 报告 3 轮 M0 评测（3 轮、覆盖全部 13 个领域，
% DeepSeek-V4-Pro、Cline v3.0.51、\texttt{thinking=low}、16 workers、2400s 超时）的各领域分解。
% 总体 Pass@1 为 \MzeroPass{}$\pm$\MzeroStd{}（3-run mean$\pm$std；r1=105/482、r2=96/486、
% r3=103/490 已判定；1500 run 中 304 PASS / 1154 FAIL / 42 ERROR-or-ABSENT）。
% 3-run 评测的逐域分解待补，将在 camera-ready 版本给出；下表保留为逐域矩阵的占位。

\begin{table}[ht]
\centering
\small
\caption{Pass@1 (\%) by domain on the 3-run M0 evaluation (\Ncases{} cases $\times$ 3 runs, DeepSeek-V4-Pro, Cline v3.0.51, \texttt{thinking=low}, 16 workers, 2400s timeout). Per-domain breakdown is pending (TODO); the Total row reports the 3-run aggregate (304 PASS / 1154 FAIL / 42 ERROR-or-ABSENT / 1500 runs; Pass\% = 304/1458 = \MzeroPass{}). Pass\% = PASS/(PASS+FAIL); ERROR and ABSENT are excluded from the denominator.}
% \caption{3 轮 M0 评测（\Ncases{} 例 $\times$ 3 轮，DeepSeek-V4-Pro、Cline v3.0.51、\texttt{thinking=low}、
% 16 workers、超时 2400s）各领域 Pass@1 (\%)。逐域分解待补（TODO）；Total 行报告 3 轮合计
%（304 PASS / 1154 FAIL / 42 ERROR-or-ABSENT / 1500 run；Pass\% = 304/1458 = \MzeroPass{}）。
% Pass\% = PASS/(PASS+FAIL)；ERROR 与 ABSENT 不计入分母。}
\label{tab:full-domain-results}
\scriptsize
\setlength{\tabcolsep}{4pt}
\begin{tabular}{@{}lcccccc@{}}
\toprule
\textbf{Domain} & \textbf{PASS} & \textbf{FAIL} & \textbf{ERROR} & \textbf{ABSENT} & \textbf{Total} & \textbf{Pass\%} \\
\midrule
schematic\_netlist & TODO & TODO & TODO & TODO & TODO & TODO \\
spice\_sim         & TODO & TODO & TODO & TODO & TODO & TODO \\
physical\_verify   & TODO & TODO & TODO & TODO & TODO & TODO \\
digital            & TODO & TODO & TODO & TODO & TODO & TODO \\
analog\_context    & TODO & TODO & TODO & TODO & TODO & TODO \\
waveform           & TODO & TODO & TODO & TODO & TODO & TODO \\
layout\_virtuoso   & TODO & TODO & TODO & TODO & TODO & TODO \\
layout\_aether     & TODO & TODO & TODO & TODO & TODO & TODO \\
verification       & TODO & TODO & TODO & TODO & TODO & TODO \\
eda\_other         & TODO & TODO & TODO & TODO & TODO & TODO \\
pg\_layout         & TODO & TODO & TODO & TODO & TODO & TODO \\
characterization   & TODO & TODO & TODO & TODO & TODO & TODO \\
pnr\_signoff       & TODO & TODO & TODO & TODO & TODO & TODO \\
\midrule
\textbf{Total}     & \textbf{304} & \textbf{1154} & \multicolumn{2}{c}{\textbf{42}} & \textbf{1500} & \textbf{\MzeroPass{}} \\
\bottomrule
\end{tabular}
\end{table}

\subsection{Per-Model Results}
\label{sec:per-model}
% \subsection{各模型结果}

This submission reports a \emph{four-backbone} comparison on the \Ncases{}-task set (3 runs each, Pass@1, same harness: Cline v3.0.51, \texttt{thinking=low}, 2400s timeout, 16 workers). DeepSeek-V4-Pro reaches \DSvPro{}$\pm$\DSvProStd{} (r1=21.78\%, r2=19.75\%, r3=21.02\%; 3/3 runs clean), doubao-seed-2.0-pro-cloud \DoubaoPro{}$\pm$\DoubaoProStd{}, doubao-seed-2.0-lite-cloud \DoubaoLite{}$\pm$\DoubaoLiteStd{}, doubao-seed-2.0-mini-cloud \DoubaoMini{}$\pm$\DoubaoMiniStd{}. A McNemar paired test on majority-vote outcomes (DeepSeek-V4-Pro vs doubao-seed-2.0-pro-cloud) gives $b{=}69$, $c{=}2$, $p\approx 0.0000$ ($\Delta{=}{+}10.67$pp, highly significant), confirming that the M0 gap is not noise. Additional backbones (GLM-5.2, Kimi-K2.6, DeepSeek-V4-Flash) are configured in the same harness but their sweeps are still running; we will report them in the camera-ready version as the multi-backbone test of RQ3 (Section~\ref{sec:per-model}).
% 本投稿给出 \Ncases{}-任务集上的\emph{四骨干}对比（各 3 轮、Pass@1、同一框架：Cline v3.0.51、
% \texttt{thinking=low}、2400s 超时、16 workers）。DeepSeek-V4-Pro 达 \DSvPro{}$\pm$\DSvProStd{}
%（r1=21.78\%、r2=19.75\%、r3=21.02\%；3/3 runs clean），doubao-seed-2.0-pro-cloud \DoubaoPro{}$\pm$\DoubaoProStd{}，
% doubao-seed-2.0-lite-cloud \DoubaoLite{}$\pm$\DoubaoLiteStd{}，doubao-seed-2.0-mini-cloud
% \DoubaoMini{}$\pm$\DoubaoMiniStd{}。对 majority-vote 结果做 McNemar 配对检验
%（DeepSeek-V4-Pro vs doubao-seed-2.0-pro-cloud）得 $b{=}69$、$c{=}2$、$p\approx 0.0000$
%（$\Delta{=}{+}10.67$pp，高度显著），印证 M0 差距非噪声。其余骨干（GLM-5.2、Kimi-K2.6、
% DeepSeek-V4-Flash）已接入同一框架但扫描仍在跑；camera-ready 版本将给出该多骨干检验作为 RQ3。

\subsection{Failure Taxonomy by Domain}
\label{sec:failure-by-domain}
% \subsection{各领域失败分类法}

{\hbadness=10000\sloppy
The batch runner records one terminal status per run, so the per-domain failure
\emph{categories} cannot be derived from that terminal status alone. We therefore
report (i) the run-level outcome anchors per domain (Table~\ref{tab:full-domain-results})
and (ii) a triage procedure that recovers the categories from each run's judge log and
agent session (\texttt{eval\_triage/05\_batch\_triage.py}, available to qualified reviewers
via the NDA-gated review package rather than publicly released, with an explicit UNKNOWN
share so that unclassified runs are never silently folded into another bucket). Applying it to
the 3-run \Ncases{}-task batch yields the per-domain category matrix in
Table~\ref{tab:failure-by-domain} (PASS runs are out of scope here; rows sum to
the FAIL+ERROR count per domain). Class abbreviations: TO=ERROR\_TIMEOUT,
INFRA=ERROR\_INFRA, NOSUB=NO\_SUBMIT, FMT=FORMAT\_CONTRACT, RT=ASSERT\_RUNTIME,
SEM=ASSERT\_SEMANTIC, ANTI=ASSERT\_ANTIHARDCODE, UNK=UNKNOWN; F+E=FAIL+ERROR.
The per-domain category counts for the 3-run batch are pending (TODO); a stacked-bar view and a Cohen's $\kappa$ agreement
statistic between the tool and a human labeler on a 100-run sample are deferred
to the camera-ready version.
\par}
% 批次 runner 每个 run 只记录一个终态，因此仅凭该终态无法推导各领域的失败\emph{类别}。
% 故我们给出：(i) 各领域 run 级结果锚点（表~\ref{tab:full-domain-results}）；
% (ii) 一个逐 run 归因流程，它从每个 run 的判分日志与智能体会话中恢复类别
% （eval_triage/05_batch_triage.py，经 NDA 受控评审包向合格评审者提供，不公开发布；
% 并显式报告 UNKNOWN 比例，避免未分类 run 被悄悄并入别的桶）。
% 将其应用于 3 轮 \Ncases{}-任务批次得到表~\ref{tab:failure-by-domain} 的逐域类别矩阵（PASS run 不在此表
% 范围；各行加和为该域 FAIL+ERROR 数）。列缩写：TO=ERROR_TIMEOUT, INFRA=ERROR_INFRA,
% NOSUB=NO_SUBMIT, FMT=FORMAT_CONTRACT, RT=ASSERT_RUNTIME, SEM=ASSERT_SEMANTIC,
% ANTI=ASSERT_ANTIHARDCODE, UNK=UNKNOWN; F+E=FAIL+ERROR。3 轮批次的逐域类别计数待补（TODO）；
% 堆叠柱状图与工具对人工标注者在 100-run 样本上的 Cohen's kappa 一致性推迟到 camera-ready。

\begin{table}[ht]
\centering
\caption{Per-domain failure-category matrix for the 3-run \Ncases{}-task batch, recovered by \texttt{eval\_triage/05\_batch\_triage.py} from per-run judge logs. PASS runs are excluded; each row sums to that domain's FAIL+ERROR count. UNKNOWN is reported as its own column, not folded into another class. Per-domain category counts are pending (TODO); overall Pass\% is \MzeroPass{} (304/1458 decided across 3 runs).}
\label{tab:failure-by-domain}
\resizebox{\columnwidth}{!}{%
\begin{tabular}{lrrrrrrrrr}
\toprule
Domain & TO & INFRA & NOSUB & FMT & RT & SEM & ANTI & UNK & F+E \\
\midrule
schematic\_netlist & TODO & TODO & TODO & TODO & TODO & TODO & TODO & TODO & TODO \\
spice\_sim         & TODO & TODO & TODO & TODO & TODO & TODO & TODO & TODO & TODO \\
physical\_verify   & TODO & TODO & TODO & TODO & TODO & TODO & TODO & TODO & TODO \\
digital            & TODO & TODO & TODO & TODO & TODO & TODO & TODO & TODO & TODO \\
analog\_context    & TODO & TODO & TODO & TODO & TODO & TODO & TODO & TODO & TODO \\
waveform           & TODO & TODO & TODO & TODO & TODO & TODO & TODO & TODO & TODO \\
layout\_virtuoso   & TODO & TODO & TODO & TODO & TODO & TODO & TODO & TODO & TODO \\
layout\_aether     & TODO & TODO & TODO & TODO & TODO & TODO & TODO & TODO & TODO \\
verification       & TODO & TODO & TODO & TODO & TODO & TODO & TODO & TODO & TODO \\
eda\_other         & TODO & TODO & TODO & TODO & TODO & TODO & TODO & TODO & TODO \\
pg\_layout         & TODO & TODO & TODO & TODO & TODO & TODO & TODO & TODO & TODO \\
characterization   & TODO & TODO & TODO & TODO & TODO & TODO & TODO & TODO & TODO \\
pnr\_signoff       & TODO & TODO & TODO & TODO & TODO & TODO & TODO & TODO & TODO \\
\midrule
\textbf{Total} & \textbf{TODO} & \textbf{TODO} & \textbf{TODO} & \textbf{TODO} & \textbf{TODO} & \textbf{TODO} & \textbf{TODO} & \textbf{TODO} & \textbf{TODO} \\
\bottomrule
\end{tabular}}
\end{table}

\subsection{Token Usage and Cost Analysis}
\label{sec:token-cost}
% \subsection{Token 用量与成本分析}

We can report \emph{behavioural} cost from the archived pilot sessions. For the 30
runs that carry an authoritative \texttt{run\_result.aggregateUsage} record, the
median run consumes 0.71M input and 55.6k output tokens (mean 0.86M/52.5k; totals
25.7M/1.58M), on top of a median of 26.5 tool calls (IQR 13.8--38.0) over 20.5
iterations and 1{,}047\,s of wall-clock time (IQR 552--1{,}495). The input:output
ratio of ${\sim}13{:}1$ is the signature of a long-horizon agent loop that re-sends
its context each iteration, and it is the main cost driver of flow orchestration.
A per-task token/cost table for the 3-run \Ncases{}-task batch is pending because the
batch summary records only terminal status and wall time; the per-run session logs
are retained and will be aggregated with the same parser
(\texttt{eval\_triage/trace\_anatomy.py}) in the camera-ready version.
% 我们能报告的是\emph{行为学}成本。对 30 份带有权威 run_result.aggregateUsage 记录的 run：
% 中位每 run 消耗 0.71M input / 55.6k output token（均值 0.86M / 52.5k；合计 25.7M / 1.58M），
% 同时中位发出 26.5 次工具调用（IQR 13.8--38.0）、跨 20.5 轮迭代、耗 1,047 秒墙钟（IQR 552--1,495）。
% 约 13:1 的 input:output 比是长时程智能体回路的特征——每轮迭代都要重发上下文，这也是流程编排的主要成本来源。
% 3 轮 \Ncases{}-任务批次的逐任务 token/成本表待补：批次摘要只记录终态与墙钟，逐 run 会话日志已留存，
% 将在 camera-ready 用同一解析器（eval_triage/trace_anatomy.py）聚合给出。

\subsection{Context-Compaction Ablation}
\label{sec:compaction-ablation}
% \subsection{上下文压缩消融}

Long-horizon agent loops accumulate context until the coding agent's compaction
mechanism must summarize earlier turns to fit the context window. Cline offers
two compaction strategies: \emph{basic} (a structured summarization that
preserves tool-call I/O verbatim) and \emph{agentic} (an LLM-generated free-form
summary). Because compaction is triggered deep into a run and its effect is
entangled with task difficulty, we isolate it with an ABBA-designed ablation on
a fixed 40-task subset of the \Ncases{}-task benchmark, holding everything else constant
(DeepSeek-V4-Pro, Cline v3.0.51, \texttt{thinking=medium}, 2400s timeout, single
trace per task). The ABBA order (basic, agentic, agentic, basic, basic, agentic)
balances position effects across six rounds.

\begin{table}[ht]
\centering
\small
\caption{Context-compaction ablation (ABBA, 6 rounds, 40-task subset). Pass\% is
PASS/40. Basic compaction consistently outperforms agentic by ${\sim}10$pp while
keeping wall-clock time lower, suggesting that preserving tool-call I/O verbatim
is more valuable than a lossy semantic summary for EDA flow orchestration.}
\label{tab:compaction-ablation}
\footnotesize
\setlength{\tabcolsep}{4pt}
\begin{tabular}{@{}llcccc@{}}
\toprule
\textbf{Round} & \textbf{Variant} & \textbf{Pass\%} & \textbf{PASS/40} & \textbf{Wall (h)} & \textbf{Tok/task (M)} \\
\midrule
r1 & basic   & 67.5 & 27 & 1.54 & 2.59 \\
r1 & agentic & 57.5 & 23 & 2.39 & 1.90 \\
r2 & agentic & 60.0 & 24 & 1.88 & 2.03 \\
r2 & basic   & 70.0 & 28 & 1.83 & 2.06 \\
r3 & basic   & 70.0 & 28 & 1.80 & 2.03 \\
r3 & agentic & 60.0 & 24 & 1.93 & 1.82 \\
\midrule
\multicolumn{2}{l}{basic mean}   & 69.2 & 27.7 & 1.72 & 2.23 \\
\multicolumn{2}{l}{agentic mean} & 59.2 & 23.7 & 2.07 & 1.92 \\
\bottomrule
\end{tabular}
\end{table}

Across the three basic rounds the pass rate is 67.5/70.0/70.0 (mean 69.2\%) and
across the three agentic rounds 57.5/60.0/60.0 (mean 59.2\%), a stable
${\sim}10$pp gap in favour of basic compaction. The ABBA design rules out a
pure ordering artifact, and the gap is consistent in direction across all three
round pairs. Interestingly, basic compaction is also \emph{faster} (1.72\,h vs
2.07\,h mean batch wall time) despite consuming slightly more tokens (2.23\,M vs
1.92\,M per task): the agentic summarizer spends extra iterations re-deriving
information that the basic summarizer keeps verbatim. We interpret this as
evidence that in EDA flow orchestration the high-value context is tool-call I/O
(exact commands, exact tool responses, exact file paths), which a free-form
semantic summary tends to lose. The ablation uses \texttt{thinking=medium}
rather than the \texttt{thinking=low} of the main evaluation, so the absolute
pass rates are not directly comparable to Table~\ref{tab:full-domain-results};
the relative comparison between the two compaction strategies is the point.
% 长时程智能体回路不断累积上下文，直到编程智能体的压缩机制必须摘要早期轮次以适配上下文窗口。
% Cline 提供两种压缩策略：\emph{basic}（保留工具调用 I/O 原文的结构化摘要）与 \emph{agentic}
% （LLM 生成的自由式摘要）。由于压缩在 run 深处触发且其效果与任务难度耦合，我们用 ABBA 设计的消融
% 在 \Ncases{}-任务基准的固定 40 任务子集上隔离该因素，其余保持恒定（DeepSeek-V4-Pro、Cline v3.0.51、
% \texttt{thinking=medium}、2400s 超时、每任务单轨迹）。ABBA 顺序（basic, agentic, agentic, basic,
% basic, agentic）在六轮间平衡位置效应。三轮 basic 的通过率为 67.5/70.0/70.0（均值 69.2\%），
% 三轮 agentic 为 57.5/60.0/60.0（均值 59.2\%），basic 稳定优 ${\sim}10$pp。ABBA 设计排除了纯顺序假象，
% 且三轮对的方向一致。有趣的是 basic 同时更快（均值 1.72h vs 2.07h）但略多耗 token（2.23M vs 1.92M/任务）：
% agentic 摘要器需额外迭代重新推导 basic 摘要器保留为原文的信息。我们将此解读为：在 EDA 流程编排中，
% 高价值上下文是工具调用 I/O（精确命令、精确工具响应、精确文件路径），自由式语义摘要倾向于丢失这些。
% 该消融使用 \texttt{thinking=medium} 而非主评测的 \texttt{thinking=low}，故绝对通过率与
% 表~\ref{tab:full-domain-results} 不可直接比较；两种压缩策略的相对比较才是重点。
